% !TeX TS-program = xelatex

\documentclass{resume}

\usepackage{graphicx}


% 更改项目符号
\renewcommand{\labelitemi}{$\bullet$}
\renewcommand{\labelitemii}{$\star$}
\renewcommand{\labelitemii}{$\ast$}

\ResumeName{马祥芸}


% 如果想插入照片，请使用以下两个库。
% \usepackage{graphicx}
% \usepackage{tikz}

\begin{document}

\ResumeContacts{
  (+86)177-250-25618,%
  \ResumeUrl{WeChatID}{WeChatID:Richardxm2062},%
  \ResumeUrl{Mail}{Email:richardxm2062@outlook.com},%
}

% 如果想插入照片，请取消此代码的注释。
% 但是默认不推荐插入照片，因为这不是简历的重点。
% 如果默认的照片插入格式不能满足你的需求，你可以尝试调整照片的大小，或者使用其他的插入照片的方法。
% 不然，也可以先渲染 PDF 简历，然后用其他工具在 PDF 上叠加照片。
% \begin{tikzpicture}[remember picture, overlay]
%   \node [anchor=north east, inner sep=1cm]  at (current page.north east) 
%      {\includegraphics[width=2cm]{image.png}};
% \end{tikzpicture}

\ResumeTitle

\section{教育经历}
\ResumeItem
[重庆大学|本科生]
{重庆大学}
[\textnormal{理论物理，物理学院，理学学士，GPA: 3.25/4.0\small{(专业前10/31)(年级前25/106)}}｜ ]   [2016.09—2020.06] 


\section{获奖}
\begin{itemize}
  \item 2017年:丙等奖学金;优秀共青团员;科技文化展一等奖（校级）
  \item 2018年:重庆大学物理学术竞赛一等奖(校级);西南地区大学生物理学术竞赛一等奖(国家级);中国大学生物理学术竞赛三等(国家级);卓越杯物理实验竞赛二等(国家级)
  \item 2019年:丙等奖学金
\end{itemize}


\section{技术学习}
\begin{itemize}
  \item \textbf{编程语言}:
  \begin{itemize}
    \item 课程学习:\textbf{[}\textbf{C};\textbf{Perl};\textbf{Fortran}\textbf{]}
    \item 自主学习:\textbf{[}\textbf{Python};\textbf{Html};文本排版语言\textbf{LaTex};符号计算\textbf{Mathematica}\textbf{]}

  \end{itemize}
  

  \item \textbf{编程基础}:
  \\ $\star$ 数据结构:熟悉数组、链表、栈、队列、哈希表、树(数组、链表构成)、堆、图(数组与链表构成)。
  \\ $\star$ 算法知识:
    \begin{itemize}
      \item 排序:冒泡排序,快速排序,计数排序,桶排序,堆排序
      \item 搜索:广度优先(bfs),深度优先(dfs)
      \item 范式:分治,回溯,动态规划,贪心算法
    \end{itemize}
  
  \item \textbf{其他}: 
  \\ $\star$ 三维建模\textbf{Creo};物理场仿真\textbf{Comsol};3D打印
\end{itemize}


\section{在校经历}
\begin{itemize}
  \item \textbf{2016年}:加入\textbf{物理学院科协学生会},积极组织参加院学生会活动;参加\textbf{重庆大学知识竞赛}并获得\textbf{团体三等奖}
  \item \textbf{2017年}:成为\textbf{班长}管理班级;参加\textbf{校青骨班素质拓展活动};加入由刘院长创立的\textbf{智能自组装机器人实验室}(半年);参加数学学院举办的\textbf{数模竞赛};初步接触全国大学生物理学术竞赛;院足球队
  \item \textbf{2018年}:参加\textbf{中国大学生物理学术竞赛}(校级-->国家级);选拔面试下一届竞赛成员;协助建立新的竞赛实验室;参加\textbf{3D打印社团}(校一级社团),管理杭州先临三维(赞助商)创客圈网站重大板块与3d打印机，并组织举办了相关征贴活动。
  \item \textbf{2019年}:参加\textbf{卓越杯物理实验竞赛}
  \item \textbf{2020年}:\textbf{毕设优秀评级}(前10\%);加入\textbf{低温物理实验室},毕业后成为科研助理,创建实验室html主页
\end{itemize}


\section{个人总结}

\begin{itemize}
  \item 
  我是一名怀有好奇心，热爱学习和不断成长的人。
  在大学期间，我取得了优异的学习成绩，同时积极加入实验室、参加各种学术竞赛，丰富了知识和技能，具有较强的自我学习和管理能力。
  同时我勇于尝试学习之外的事物，成为班长管理班级、加入院足球队、参加院学生会、参加校社团等。锻炼了我的组织和领导能力，不断拓展自己的兴趣和社交能力，具有良好的沟通能力的团队合作精神。
  我善于在压力下保持高效工作，并自我驱动学习新的事物和技能以解决问题，能够迅速适应新的工作环境，并且始终保持着对知识的渴望和追求。
  尽管我目前在职业领域中没有过多的实践经验，但我有信心通过自己的努力和学习，迅速适应工作，并且胜任其中。
  我对新事物充满好奇心，愿意接受挑战，不断突破自己的局限，实现个人和团队的共同目标。
  
  
  
\end{itemize}

\end{document}
